\chapter{Concluzii}
\label{cap:cap5}

Lucrarea a urmărit dezvoltarea unei platforme web pentru managementul
fluxurilor studențești, cu accent pe monitorizarea prezenței la laboratoare.
Obiectivele stabilite în introducere au fost atinse prin definirea unei
arhitecturi client-server, implementarea unui frontend modular în Angular,
realizarea unui backend REST în PHP și proiectarea unei baze de date relaționale
în MySQL.

Contribuțiile principale sunt:

\begin{itemize}
    \item integrarea într-o singură aplicație a fluxurilor pentru student,
          profesor și administrator;
    \item automatizarea procesului de marcare și consultare a prezențelor;
    \item organizarea modulară a codului, care permite extindere ulterioară;
    \item definirea unui model de date coerent, cu relații și constrângeri explicite.
\end{itemize}

Rezultatele obținute indică faptul că soluția reduce efortul administrativ,
oferă vizibilitate mai bună asupra activității didactice și susține o
gestionare mai transparentă a datelor academice.

\section{Direcții viitoare de dezvoltare}
\label{cap:cap5:directii-viitoare}

Etapele următoare de dezvoltare vizează:

\begin{itemize}
    \item introducerea testării automate extinse (unitare și end-to-end);
    \item optimizări de performanță pentru volume mari de date;
    \item extinderea modelului pentru discipline opționale și scenarii complexe;
    \item consolidarea componentelor de securitate și audit al operațiilor.
\end{itemize}

Per ansamblu, platforma reprezintă o bază solidă pentru digitalizarea
proceselor de management academic și poate fi adaptată la cerințe instituționale
mai ample.